% S16 Appendix. Convolution operators and controls for the tissue domains, in full.
% PLOS ONE PONE-D-26-04349, Osgood-Zimmerman & Raredon. Supporting Information.
% Text moved verbatim from the main text of the 2026-10-08 revision when the paper was reframed
% around uncertainty (2026-10-22); only cross-references were adapted.
% Compile: pdflatex -> bibtex -> pdflatex x2.
\documentclass[10pt]{article}
\usepackage[margin=1in]{geometry}
\usepackage{amsmath,amssymb}
\usepackage{cite}
\usepackage[hidelinks]{hyperref}
\usepackage{parskip}
\bibliographystyle{plos2025}
\begin{document}
\begin{center}
{\large\bfseries S16 Appendix. Convolution operators and controls for the tissue domains.}\\[4pt]
\emph{Osgood-Zimmerman \& Raredon}
\end{center}
\section*{Comparison of convolution operators}


We assessed the output spatial structure that resulted from clustering 32 distinct spatial assays, summarized in S5 File. Three in particular are detailed in S15 Fig. Of the operators we tested, the `common minimum' approach, which simply takes the common minimum between the ligand and receptor subunits, gave the largest number of clusters (29, against 27 for both the product and the geometric mean). Its spatial entropy was close to the lowest of the operators and well below that of the product: mean ELSA at radius 100 was 0.050 for the common minimum, 0.047 for the geometric mean, and 0.071 for the product. Of note, the common minimum approach closely approximates the outputs produced when $k$, the dissociation constant, is decreased to 0, which within that kinetic model corresponds to total binding without any dissociation. Because of these results, we carried this output forward for downstream experimentation. For full results from this experiment, see S6 File.

Two caveats attach to this ranking. It rests on single runs of a clustering step whose run-to-run variation we have not measured directly; re-runs of the related clustering in the domain figure of the main text differed from the published values by up to about 0.01 ELSA units (see the controls below). At radius 50 the common minimum and the product differ by 0.024, about twice that, while the common minimum and the geometric mean differ by 0.004, well within it, so those two should not be read as ranked on coherence. It is also a comparison of operators within our own pipeline on one dataset, not a general result about which convolution is correct.


\section*{Controls for the tissue domains}


The experiments above show that clustering the full signaling hyperfield yields spatially coherent domains, but not that those domains reflect structure present in the tissue. Two properties of the pipeline could in principle manufacture coherent domains on their own: the spatial model smooths, and the clustering operates in a space of more than a thousand dimensions. We therefore ran two controls aimed at these two properties, and a third, described after them, aimed at the ligand-receptor pairing. In each, the domain-generating pipeline is held fixed between the observed and the control runs, so that any difference is attributable to the manipulation alone. Spatial coherence is quantified throughout with ELSA at neighborhood radii of 50 and 100 units (see the Methods section); ELSA is negatively oriented, so lower values indicate more spatially coherent domains. The first two controls are summarized in S10 Fig.

\subsection*{A coordinate shuffle destroys the domains}

To test whether the spatial smoother manufactures coherent structure from data carrying no spatial signal, we permuted the assignment of the observed counts to spatial locations. This destroys spatial structure while preserving the observed count distribution exactly. We then refit the spatial model for every one of the 912 features from scratch on the permuted data, and passed the refit fields through the identical convolution, clustering, and scoring pipeline. Identical here means identical between the observed and the permuted data, which is what the comparison requires. It is not identical to the clustering behind the domain figure of the main text: because each permutation already required 912 model fits, both arms of this control were clustered with a faster implementation of the same two steps, a nearest-neighbor graph built with RcppHNSW and Louvain clustering in igraph, in place of Seurat. ELSA values from the two implementations are not on the same scale, so the values in this section should be compared with each other and not with those reported for the domain figure of the main text. Refitting, rather than reshuffling fields that were fitted on the real coordinates, is what makes this a test of the smoother: the model sees only spatially meaningless data. Because each permutation requires 912 model fits, the experiment was run on a high-throughput computing cluster, with 49 permutations.

The observed data yielded 20 domains, with an ELSA at radius 50 of 0.075. The 49 coordinate-shuffled replicates yielded between 4 and 12 domains, with ELSA at radius 50 averaging 0.470 and ranging from 0.412 to 0.583. The null mean is 6.3 times the observed value, and the lowest of the 49 shuffled replicates is still 5.5 times it, so the observed and null values do not overlap. The permutation $p$-value is 0.020, which is the smallest value attainable with 49 permutations; the test is therefore limited by the number of replicates we could afford rather than by any overlap between the observed value and the null distribution.

Fitting a spatial Gaussian process to spatially permuted data does not, then, produce coherent tissue domains: it produces a small number of incoherent fragments. Because the permuted data pass through a pipeline of identical dimensionality, this speaks to the smoothing and the dimensionality explanations at once.

\subsection*{The domains correspond to independently annotated anatomy}

As a positive control, we compared the emergent domains against the anatomical annotations published with the source dataset \cite{Chen2022} for the same tissue section. These annotations were produced by a different group, by a different method, and without any reference to our labels, which our pipeline in turn never sees; in that sense they are an independent reference. All 26,000 modeled bins matched an annotated bin.

For completeness about what the reference is: Chen et al.\ derived these labels by spatially constrained clustering, combining a nearest-neighbor graph over spatial position with one over expression, and then labeling each resulting cluster by hand from its marker genes against a reference embryo atlas. The comparison is therefore between two partitions of the same tissue section, which is unavoidable, since a concordance check is only meaningful against annotations of the same specimen. The two partitions are nonetheless built from very different features: theirs from binned expression across all genes, ours from the inferred interaction fields of 1477 ligand-receptor pairs, and the manual labeling step brings in external anatomical knowledge from outside the dataset entirely.

Agreement was well above chance. The adjusted Rand index between our domains and the anatomical labels was 0.401 and the normalized mutual information was 0.604, against a permutation null whose mean adjusted Rand index was 0.00003 and whose maximum over 499 permutations was 0.001, giving $p = 0.002$, again the floor for the number of permutations used. The mean size-weighted purity of the emergent domains with respect to the anatomical labels was 77.6\%. Several individual domains recover recognized structures cleanly: spinal cord at 97\% purity, liver at 96\%, one heart domain at 99\%, and meninges at 92\%, with lung, muscle, thymus, and cartilage primordium also recovered.

We would not read an adjusted Rand index of 0.401 as showing that the signaling domains reproduce the anatomical segmentation; it is well above chance but far from perfect agreement. The disagreement is also structured rather than random, in that our domains subdivide several organs, most visibly the heart, into multiple territories. We interpret this as a consequence of partitioning the tissue on a different quantity: the published annotation partitions by overall expression profile, labeled by cell-type markers, whereas these domains partition by inferred ligand-receptor signaling context, and there is no reason for the two to coincide. The conclusion we draw is the narrower one, that the emergent domains recover major anatomical structures while resolving additional substructure within them.

The two controls answer different questions, and we do not want them read as interchangeable. The coordinate shuffle asks whether the structure is real, and it is the one that can rule out the smoothing and dimensionality explanations, because it manipulates the data and refits everything. The concordance comparison asks whether the structure that survives is anatomically meaningful, which is what makes it the natural complement rather than a substitute.

\subsection*{Reproducibility of the clustering}

Two implementations of the clustering step appear in this work, and their reproducibility has to be stated separately. S15 Fig and the domain figure of the main text use Seurat. The coordinate-shuffle control, and one arm of the pair-shuffling control below, use a faster implementation of the same two steps, a nearest-neighbor graph built with RcppHNSW followed by Louvain clustering in igraph, which made it affordable to cluster many permuted datasets. The two implementations give ELSA values on different scales, so each is compared only with itself.

The faster implementation, as we first configured it, did not reproduce exactly from run to run. Its nearest-neighbor graph was built across four threads without a fixed random seed, and under concurrent insertion the graph depends on the order in which threads happen to complete. Across eight replicate runs on identical input, ELSA at radius 50 ranged from 0.053 to 0.067, with a standard deviation of 0.0045, and the number of domains ranged from 19 to 22. Two conditions together make it exactly reproducible, and neither suffices alone. With four threads and a fixed seed, three replicate runs still gave 19, 19 and 20 domains; unseeded single-threaded runs gave 19, 22 and 22. With a single thread \textit{and} a fixed seed, five replicate runs returned bit-identical cluster assignments, with 21 domains in all five, and a second seed likewise gave three identical runs, at 23 domains. The seed therefore changes the result, and is a parameter of the analysis to be reported rather than an incidental detail. The coordinate-shuffle control was run before this was established; its effect is far larger than this run-to-run spread, so its conclusion does not depend on it.

For the Seurat clustering behind S15 Fig and the domain figure of the main text we have not established whether the result is exactly reproducible, and we have not measured its run-to-run variation directly. As an indirect check, an independently written script re-ran the variance-ordered sweep at six depths (10, 50, 150, 350, 650 and 1000 fields) and the full-depth clustering. Mean ELSA at radius 50 was within 0.012 of the published value at each of the six depths, and at full depth it was 0.1236 against the published 0.1237, with 19 domains against the published 20. These re-runs omitted the four duplicated input columns described in the Materials and methods section, so they differ from the published runs in input as well as in run. We therefore treat each domain map in S15 Fig and the domain figure of the main text as a single realization, read its domain count as approximate, and regard differences in ELSA of about 0.01 between single Seurat runs as within what a re-run could produce.

\subsection*{Shuffling the ligand-receptor pairing does not make the domains less coherent}

A third control asks whether the domains depend on which receptor is paired with which ligand. We reassigned receptors to ligands at random, then repaired by targeted swaps every reassigned pair that coincided with a real FANTOM5 pairing or paired a gene with itself, so that no shuffled set retained a real pair. The set of receptors is unchanged, so every shuffled set is built from exactly the same fitted fields as the true pairing, and only the pairing differs. Each shuffled set then passed through the same convolution, clustering and scoring as the true pairing.

In Seurat, with 19 shuffled sets clustered at full depth, the true pairing gave an ELSA at radius 50 of 0.124 against a shuffled mean of 0.093 (standard deviation 0.018), so 17 of the 19 shuffles produced more coherent domains than the true pairing did; the true pairing gave 19 domains against a shuffled mean of 18.5. In the faster implementation, run with a single thread and a fixed seed on 99 shuffled sets, the true pairing sat inside the shuffled distribution: 24 of the 99 shuffles were more coherent, and it gave 21 domains against a shuffled mean of 21.05. The direction of the small difference is not consistent between the two implementations, and we read the result as no detectable effect of the pairing on the coherence of the domains.

The pairing is not without consequence, however. True pairings yield 1,449 numerically distinct fields. Because this count needs no clustering, we computed it for 9,999 shuffled pairings, made with the same procedure and including the 99 clustered ones: they gave a mean of 1,424 (standard deviation 5.9), and only one reached 1,449 or more, a permutation $p$ of 0.0002. True partners are also more alike in spatial pattern than random partners. For each of the 1,477 pairs we compared the correlation between the two true partner fields with the correlation between the ligand field and that of a different receptor that is not a FANTOM5 partner of the ligand, chosen from 30 random candidates to match the true pair as closely as possible in how consistently one partner is the smaller of the two. The mean Pearson correlation was 0.111 for true partners and 0.068 for random ones, true partners were the more correlated in 60.8\% of pairs, and a paired $t$-test gave $p = 4 \times 10^{-21}$. Even random partners are positively correlated (mean Spearman correlation 0.325), consistent with a spatial pattern shared by most of these genes, so the part of the correlation specific to the pairing is modest. We return to what these results imply in the Discussion of the main text.


\bibliography{references-aoz,references-msbr}
\end{document}
